\NSPage{145}%
vanish and \NSi{NS-F-P145-001}{\Pi_0} is holomorphic. For
\NSi{NS-F-P145-002}{\Lambda\geq 1} define
\NSFormula{NS-F-P145-003}{display}{B.3}
\begin{equation}
  \zeta_*=-\frac{LH_*}{H_*^2+\sigma_*^2},
  \qquad \widetilde\xi_0=\Lambda\zeta_*,
  \qquad \phi_*=\exp\!\left(\Lambda\int_0^\eta \zeta_*(w)\,dw\right).
  \tag{B.3}\label{eq:B.3}
\end{equation}
The integral is well defined on \NSi{NS-F-P145-004}{\Omega}. On the real interval
\NSi{NS-F-P145-005}{\phi_*>0}. The amplitude \NSi{NS-F-P145-006}{\mathsf C} will be chosen after
\NSi{NS-F-P145-007}{\Lambda}.

\subsection{An analytic coefficient space.}\label{sec:B.2}
The profile equations \eqref{eq:4.13} contain parameter derivatives as well as radial derivatives,
with a regular singular point at \NSi{NS-F-P145-008}{X=0}. In this subsection we introduce a
weighted coefficient space \NSi{NS-F-P145-009}{\mathcal B_\rho} with norm \eqref{eq:B.4} for their
integral form, in which increasing the radial degree compensates for a parameter derivative. The
bounds in Lemma~\ref{lem:B.1} control multiplication and radial inversion in this space, including
the mixed nonlinear terms needed in the contraction argument for Proposition~\ref{prop:B.2}.

For the analytic axis construction, we use the scalar radial variable
\NSi{NS-F-P145-010}{Y=\Lambda X}. Fix a small radius \NSi{NS-F-P145-011}{\rho>0}. For
\NSi{NS-F-P145-012}{F(Y,\eta)=\sum_{\alpha\geq0}F_\alpha(\eta)Y^\alpha} put
\NSFormula{NS-F-P145-013}{display}{B.4}
\begin{equation}
  \mathfrak a_{\alpha\beta}
  =\frac{20^{-\alpha}\rho^{-\beta}\beta!\binom{\alpha+\beta}{\beta}}
         {(\alpha+1)^2(\beta+1)^2},
  \qquad
  \lVert F\rVert_\rho
  =\sup_{\alpha,\beta\geq0}\sup_{\eta\in I}
    \frac{|\partial_\eta^\beta F_\alpha(\eta)|}{\mathfrak a_{\alpha\beta}}.
  \tag{B.4}\label{eq:B.4}
\end{equation}
Let \NSi{NS-F-P145-014}{\mathcal B_\rho} consist of the sequences of smooth coefficients with
finite norm. It is complete: a Cauchy sequence converges uniformly at each coefficient and
derivative order, and uniform convergence of successive derivatives identifies those limits as
the derivatives of the limiting coefficient. In the rescaled variable \NSi{NS-F-P145-015}{Y},
the logarithmic radial derivative is
\NSi{NS-F-P145-016}{\mathcal D_X=X\partial_X=Y\partial_Y}. Let
\NSi{NS-F-P145-017}{\mathcal I F=\int_0^Y F(Y',\eta)\,dY'}. Radial averaging is unchanged by
the scaling \NSi{NS-F-P145-018}{Y=\Lambda X}, so
\NSi{NS-F-P145-019}{\mathcal A_X(F)=Y^{-1}\mathcal I F}.

\begin{lemma}\label{lem:B.1}
Multiplication is bounded on \NSi{NS-F-P145-020}{\mathcal B_\rho}. For
\NSi{NS-F-P145-021}{\nu=1,2}, let \NSi{NS-F-P145-022}{\mathcal J_\nu F} be the solution
\NSi{NS-F-P145-023}{G} of
\NSi{NS-F-P145-024}{YG_{YY}+\nu G_Y=F} that extends smoothly to
\NSi{NS-F-P145-025}{Y=0} and satisfies \NSi{NS-F-P145-026}{G(0)=0}. Its coefficients are given by
\NSFormula{NS-F-P145-027}{display}{B.5}
\begin{equation}
  (\mathcal J_\nu F)_{\alpha+1}
  =\frac{F_\alpha}{(\alpha+1)(\alpha+\nu)},
  \qquad (\mathcal J_\nu F)_0=0.
  \tag{B.5}\label{eq:B.5}
\end{equation}
The operator \NSi{NS-F-P145-028}{\mathcal J_\nu} is bounded on
\NSi{NS-F-P145-029}{\mathcal B_\rho}, as are radial averaging, multiplication by
\NSi{NS-F-P145-030}{Y}, and the operators \NSi{NS-F-P145-031}{\mathcal I} and
\NSi{NS-F-P145-032}{\partial_\eta\mathcal I}. In addition,
\NSFormula{NS-F-P145-033}{display}{B.6}
\begin{equation}
  \bigl\lVert\mathcal J_\nu[(\partial_\eta F)(\mathcal D_XG)]\bigr\rVert_\rho
  \leq C_\rho\lVert F\rVert_\rho\lVert G\rVert_\rho.
  \tag{B.6}\label{eq:B.6}
\end{equation}
The same estimate holds if either derivative is omitted or \NSi{NS-F-P145-034}{F} is replaced by
\NSi{NS-F-P145-035}{\mathcal A_X(F)}. Extra undifferentiated factors can be inserted, with one
additional algebra constant for each factor.
\end{lemma}

\begin{proof}
The elementary convolution estimate
\NSFormula{NS-F-P145-036}{display}{B.7}
\begin{equation}
  \sum_{i=0}^{N}\frac{(N+1)^2}{(i+1)^2(N-i+1)^2}
  \leq 8\sum_{i=1}^{\infty}i^{-2}=:C_{\mathrm{sq}}
  \tag{B.7}\label{eq:B.7}
\end{equation}
follows by dividing the sum at \NSi{NS-F-P145-037}{N/2}. In the Leibniz formula for the product
coefficient, its derivative binomial cancels the factorials in the weights. The remaining binomial
ratio is at most one, since
\NSFormula{NS-F-P145-038}{display}{B.8}
\begin{equation}
  \binom{i+k}{k}\binom{\alpha-i+\beta-k}{\beta-k}
  \leq\binom{\alpha+\beta}{\beta}.
  \tag{B.8}\label{eq:B.8}
\end{equation}
Indeed the left side counts some of the \NSi{NS-F-P145-039}{\beta}-element subsets of a set split
into blocks of sizes \NSi{NS-F-P145-040}{i+k} and
\NSi{NS-F-P145-041}{\alpha-i+\beta-k}. Applying \eqref{eq:B.7} in both indices gives the algebra
constant \NSi{NS-F-P145-042}{C_{\mathrm{sq}}^2}.

\NSPage{146}%
The two weight ratios needed for integration are
\NSFormula{NS-F-P146-001}{display}{B.9}
\begin{equation}
  \frac{\mathfrak a_{\alpha\beta}}{\mathfrak a_{\alpha+1,\beta}}
  =20\frac{(\alpha+2)^2}{(\alpha+1)^2}\frac{\alpha+1}{\alpha+\beta+1}
  \leq80,
  \tag{B.9}\label{eq:B.9}
\end{equation}
\NSFormula{NS-F-P146-002}{display}{B.10}
\begin{equation}
  \frac{\mathfrak a_{i,k+1}}{\mathfrak a_{i+1,k}}
  =\frac{20}{\rho}(i+1)\frac{(i+2)^2(k+1)^2}{(i+1)^2(k+2)^2}
  \leq\frac{80}{\rho}(i+1).
  \tag{B.10}\label{eq:B.10}
\end{equation}
The first proves the asserted bounds for \NSi{NS-F-P146-003}{\mathcal J_\nu},
\NSi{NS-F-P146-004}{\mathcal I}, and multiplication by \NSi{NS-F-P146-005}{Y}; averaging
divides degree \NSi{NS-F-P146-006}{\alpha} by \NSi{NS-F-P146-007}{\alpha+1}. The second proves
the bound for \NSi{NS-F-P146-008}{\partial_\eta\mathcal I}, because its integration divisor
cancels \NSi{NS-F-P146-009}{i+1}.

For the mixed derivative bound \eqref{eq:B.6}, take the degree
\NSi{NS-F-P146-010}{N=\alpha+1} coefficient of its left side. Assign its extra degree to the factor
carrying the parameter derivative, and write \NSi{NS-F-P146-011}{I_1=i+1},
\NSi{NS-F-P146-012}{J_1=\alpha-i}. After \eqref{eq:B.10} and division by
\NSi{NS-F-P146-013}{\mathfrak a_{N\beta}}, its \NSi{NS-F-P146-014}{\beta}th derivative is bounded by
\NSFormula{NS-F-P146-015}{display}{none}
\begin{equation*}
  \frac{80}{\rho}\lVert F\rVert_\rho\lVert G\rVert_\rho
  \sum_{i=0}^{\alpha}\sum_{k=0}^{\beta}
  \frac{I_1J_1}{N(\alpha+\nu)}
  \frac{(N+1)^2}{(I_1+1)^2(J_1+1)^2}
  \frac{(\beta+1)^2}{(k+1)^2(\beta-k+1)^2}.
\end{equation*}
The binomial ratio was dropped using \eqref{eq:B.8} with radial indices
\NSi{NS-F-P146-016}{I_1,J_1}. The first fraction is at most one, and the two sums are bounded by
\eqref{eq:B.7}. With an average on \NSi{NS-F-P146-017}{F}, its divisor
\NSi{NS-F-P146-018}{i+1} cancels \NSi{NS-F-P146-019}{I_1}; the resulting fraction is again at
most one. If the logarithmic radial derivative is omitted, replace
\NSi{NS-F-P146-020}{J_1} by one. A logarithmic radial derivative without an
\NSi{NS-F-P146-021}{\eta}-derivative is controlled directly by \eqref{eq:B.9} and
\eqref{eq:B.5}. For additional factors, the same allocation leaves their indices undifferentiated
and repeated convolution gives the claim. This argument includes the endpoint indices and every
\NSi{NS-F-P146-022}{\beta}.
\end{proof}

\subsection{The nonlinear analytic axis profile.}\label{sec:B.3}
We now solve the equations \eqref{eq:4.13} with vanishing leading residual stress for the chosen
axis data (Section~\ref{sec:B.1}). After rescaling the radius by
\NSi{NS-F-P146-023}{Y=\Lambda X}, we will show that the solution is a small perturbation of an
explicit profile as \NSi{NS-F-P146-024}{\Lambda} grows (\eqref{eq:B.13}). The scalar comparison
function \NSi{NS-F-P146-025}{f_0} below controls the normalized azimuthal profile
\NSi{NS-F-P146-026}{\phi/\phi_*} through \NSi{NS-F-P146-027}{f_0(Y\chi)}; we will verify that it
remains positive throughout the required radial interval in \eqref{eq:B.11}. This will ensure that
division by \NSi{NS-F-P146-028}{\phi} in the profile equations is valid.

The analytic scalar comparison function is
\NSFormula{NS-F-P146-029}{display}{none}
\begin{equation*}
  f_0(z)=\sum_{\alpha\geq0}\frac{(-z/2)^\alpha}{\alpha!(\alpha+1)!}.
\end{equation*}
For \NSi{NS-F-P146-030}{0\leq z\leq4.1}, set \NSi{NS-F-P146-031}{t=z/2}. The alternating tail
after the cubic term has decreasing magnitudes, and hence
\NSFormula{NS-F-P146-032}{display}{B.11}
\begin{equation}
  f_0(z)\geq1-\frac t2+\frac{t^2}{12}-\frac{t^3}{144}
  \geq\frac{305719}{1152000}>.265,
  \qquad f_0(z)\leq1.
  \tag{B.11}\label{eq:B.11}
\end{equation}
The cubic is decreasing on \NSi{NS-F-P146-033}{[0,2.05]}, giving its stated endpoint value. For
the upper bound group the alternating tail after the constant term; its magnitudes decrease from
the linear term onward.

\begin{proposition}\label{prop:B.2}
With the axis data of Section~\ref{sec:B.1}, there are \NSi{NS-F-P146-034}{\Lambda_0} and, for
each \NSi{NS-F-P146-035}{\Lambda\geq\Lambda_0}, an amplitude threshold
\NSi{NS-F-P146-036}{\mathsf C_0(\Lambda)} such that every
\NSi{NS-F-P146-037}{\mathsf C\geq\mathsf C_0(\Lambda)} admits an analytic axis profile with vanishing leading
residual stress on \NSi{NS-F-P146-038}{0\leq Y=\Lambda X\leq4.1}. It has the form
\NSFormula{NS-F-P146-039}{display}{B.12}
\begin{equation}
  \phi=\phi_*\Phi,
  \qquad U=U_*+\Lambda^{-1}\mathfrak u,
  \qquad \Pi=\Pi_0+\Lambda^{-1}\mathcal I(g^2\Phi^2),
  \qquad g=\phi_*/\mathsf C,
  \tag{B.12}\label{eq:B.12}
\end{equation}
\NSPage{147}%
where \NSi{NS-F-P147-001}{\Phi(0,\eta)=1} and \NSi{NS-F-P147-002}{\mathfrak u(0,\eta)=0}. The profile
is analytic in \NSi{NS-F-P147-003}{Y} and in \NSi{NS-F-P147-004}{\eta} on a common
neighborhood of \NSi{NS-F-P147-005}{[0,4.1]\times I}. For every fixed
\NSi{NS-F-P147-006}{r,s\geq0},
\NSFormula{NS-F-P147-007}{display}{B.13}
\begin{equation}
  \sup_{[0,4.1]\times I}
  \left|\partial_Y^r\partial_\eta^s
    \left(\Phi-f_0(Y\chi),\mathfrak u+\frac{YZ_*}{2L}\right)\right|
  \leq\frac{C_{r,s}}{\Lambda}.
  \tag{B.13}\label{eq:B.13}
\end{equation}
Here the constants and a smaller common parameter neighborhood are independent of sufficiently
large \NSi{NS-F-P147-008}{\Lambda} and of \NSi{NS-F-P147-009}{\mathsf C\geq\mathsf C_0(\Lambda)}. In unscaled
radial derivatives, the corresponding bound is \NSi{NS-F-P147-010}{C_{r,s}\Lambda^{r-1}}.
\end{proposition}

\begin{proof}
We rewrite the equations in the rescaled radius so that the nonlinear remainders carry a factor
\NSi{NS-F-P147-011}{\Lambda^{-1}}. The radial inverses then reduce the problem to a contraction
about the explicit comparison profile. In the notation of the stress equations, put
\NSFormula{NS-F-P147-012}{display}{B.14}
\begin{equation}
  B=-2D_\eta\mathcal A_X(\mathfrak u)-d\partial_\eta\mathcal A_X(\mathfrak u),
  \qquad W=W_*+\Lambda^{-1}B,
  \qquad H_c=H_*+\Lambda^{-1}d\mathfrak u,
  \qquad \mathfrak p=\mathcal I(g^2\Phi^2).
  \tag{B.14}\label{eq:B.14}
\end{equation}
The equations with vanishing leading residual stress \eqref{eq:4.13}, together with radial pressure
balance, are
\NSFormula{NS-F-P147-013}{display}{B.15}
\begin{equation}
  -2L(X\phi_{XX}+2\phi_X)/\phi=S_q,
  \qquad -2L(XU_{XX}+U_X)=S_n,
  \qquad \Pi_X=\phi^2/\mathsf C^2.
  \tag{B.15}\label{eq:B.15}
\end{equation}
Substitution of \eqref{eq:B.12} gives
\NSFormula{NS-F-P147-014}{display}{none}
\begin{align*}
  2(Y\Phi_{YY}+2\Phi_Y)&=-\chi\Phi+\Lambda^{-1}R_1,\\
  2(Y\mathfrak u_{YY}+\mathfrak u_Y)&=-Z_*/L+\Lambda^{-1}R_2,
\end{align*}
with the explicit remainders
\NSFormula{NS-F-P147-015}{display}{none}
\begin{align*}
  LR_1={}&[W+h(1-2\eta U)+d\mathfrak u\zeta_*]\Phi+W\mathcal D_X\Phi+H_c\Phi_\eta,\\
  LR_2={}&[A(1-4\eta U_*)+dU_{*\eta}]\mathfrak u-2A_\eta\Lambda^{-1}\mathfrak u^2
            +W\mathcal D_X\mathfrak u+H_*\mathfrak u_\eta+d\Lambda^{-1}\mathfrak u\mathfrak u_\eta\\
          &\quad-4A_\eta\mathfrak p+d\partial_\eta\mathfrak p-2\eta\mathcal D_X\mathfrak p.
\end{align*}
The leading angular term has the stated sign because
\NSi{NS-F-P147-016}{H_*\widetilde\xi_0/(\Lambda L)=-\chi}. The expanded axial remainder follows by
subtracting the constant value \NSi{NS-F-P147-017}{-Z_*/L} from
\NSi{NS-F-P147-018}{L^{-1}\{A(1-2\eta U)U+W\mathcal D_XU+H_cU_\eta-4A_\eta\Pi+d\Pi_\eta-2\eta\mathcal D_X\Pi\}}
and multiplying by \NSi{NS-F-P147-019}{\Lambda}.

To apply the radial inverse bounds, we must control the differentiated products in
\NSi{NS-F-P147-020}{R_1} and \NSi{NS-F-P147-021}{R_2}. The only products containing both an
unintegrated parameter derivative and a logarithmic radial derivative are
\NSFormula{NS-F-P147-022}{display}{none}
\begin{equation*}
  (\partial_\eta\mathcal A_X(\mathfrak u))\mathcal D_X\Phi,
  \qquad (\partial_\eta\mathcal A_X(\mathfrak u))\mathcal D_X\mathfrak u.
\end{equation*}
Their derivatives act on distinct factors. Furthermore,
\NSi{NS-F-P147-023}{\partial_\eta\mathfrak p=\partial_\eta\mathcal I(g^2\Phi^2)} is bounded by
Lemma~\ref{lem:B.1}, and \NSi{NS-F-P147-024}{\mathcal D_X\mathfrak p=Yg^2\Phi^2}. Thus the maps
\NSi{NS-F-P147-025}{(\Phi,\mathfrak u)\longmapsto\mathcal J_\nu R_i} are bounded and locally Lipschitz on
fixed balls of \NSi{NS-F-P147-026}{\mathcal B_\rho^2}. This last assertion follows quantitatively
by the displayed algebra bounds; to estimate a difference, subtract each monomial by replacing its
factors one at a time.

The coefficient bounds just used are uniform in the large parameters. Indeed
\NSi{NS-F-P147-027}{\zeta_*,\chi,U_*,Z_*/L} are fixed holomorphic functions on
\NSi{NS-F-P147-028}{\Omega}. Choose a smaller neighborhood and take
\NSi{NS-F-P147-029}{\rho} strictly below its distance to the boundary of
\NSi{NS-F-P147-030}{\Omega}. After choosing \NSi{NS-F-P147-031}{\Lambda}, require
\NSFormula{NS-F-P147-032}{display}{B.16}
\begin{equation}
  \mathsf C\geq\sup_{\eta\in\Omega}|\phi_*(\eta)|.
  \tag{B.16}\label{eq:B.16}
\end{equation}
Then \NSi{NS-F-P147-033}{|g|\leq1} on \NSi{NS-F-P147-034}{\Omega}, so Cauchy's inequality
bounds its derivatives on the smaller neighborhood uniformly in both large parameters. The strict
choice of \NSi{NS-F-P147-035}{\rho} absorbs the factor
\NSi{NS-F-P147-036}{(\beta+1)^2} in \eqref{eq:B.4}.
\NSPage{148}%
It is essential here that the bound on \NSi{NS-F-P148-001}{g} holds on a complex neighborhood,
since \NSi{NS-F-P148-002}{\phi_*} itself can grow exponentially with
\NSi{NS-F-P148-003}{\Lambda}.

It remains to invert the order-one angular term without assuming that its coefficient is small. Let
\NSi{NS-F-P148-004}{M_\chi} bound multiplication by \NSi{NS-F-P148-005}{\chi} on
\NSi{NS-F-P148-006}{\mathcal B_\rho} and set
\NSi{NS-F-P148-007}{T=\mathcal J_2\chi/2}, where multiplication precedes integration. If
\NSi{NS-F-P148-008}{F} has no radial coefficient below degree \NSi{NS-F-P148-009}{b},
\eqref{eq:B.9} gives
\NSFormula{NS-F-P148-010}{display}{none}
\begin{equation*}
  \lVert\mathcal J_2F\rVert_\rho
  \leq\frac{80}{(b+1)(b+2)}\lVert F\rVert_\rho,
  \qquad
  \lVert T^k\rVert\leq\frac{(40M_\chi)^k}{k!(k+1)!}.
\end{equation*}
Multiplication by \NSi{NS-F-P148-011}{\chi} preserves the vanishing of coefficients below degree
\NSi{NS-F-P148-012}{b}, and \NSi{NS-F-P148-013}{\mathcal J_2} increases the minimum possible
nonzero degree by one. Parameter differentiation introduces no lower-degree coefficients, so the
estimate includes arbitrarily high parameter derivatives of the iterated products. The series
\NSi{NS-F-P148-014}{\sum_{k\geq0}(-T)^k} therefore converges absolutely in operator norm and is a
two-sided inverse of \NSi{NS-F-P148-015}{1+T}.

The unperturbed solution is
\NSi{NS-F-P148-016}{\Phi^0=(1+T)^{-1}1=f_0(Y\chi)} and
\NSi{NS-F-P148-017}{\mathfrak u^0=-YZ_*/(2L)}. Solve the equations by the map
\NSFormula{NS-F-P148-018}{display}{none}
\begin{equation*}
  (\Phi,\mathfrak u)\longmapsto
  \left(\Phi^0+\frac{(1+T)^{-1}\mathcal J_2R_1}{2\Lambda},
        \mathfrak u^0+\frac{\mathcal J_1R_2}{2\Lambda}\right).
\end{equation*}
On a fixed ball about \NSi{NS-F-P148-019}{(\Phi^0,\mathfrak u^0)} the uniform bounds and Lipschitz
constants established above make this ball invariant under the map and make the map a strict
contraction for all sufficiently large \NSi{NS-F-P148-020}{\Lambda}. Its fixed point has norm
distance at most \NSi{NS-F-P148-021}{C/\Lambda} from the center. Choosing
\NSi{NS-F-P148-022}{\mathsf C_0(\Lambda)} to include \eqref{eq:B.16} implements the stated order of
parameters.

Finally, for \NSi{NS-F-P148-023}{R<20},
\NSFormula{NS-F-P148-024}{display}{none}
\begin{equation*}
  \sum_{\alpha\geq0}\binom{\alpha+\beta}{\beta}(R/20)^\alpha
  =(1-R/20)^{-\beta-1}.
\end{equation*}
Choose \NSi{NS-F-P148-025}{4.1<R<20}. The coefficient norm therefore gives a common positive
parameter radius, any sufficiently small number below
\NSi{NS-F-P148-026}{\rho(1-R/20)}, on \NSi{NS-F-P148-027}{|Y|\leq R}. Cauchy estimates on a
smaller \NSi{NS-F-P148-028}{Y} disk give \eqref{eq:B.13}, including all fixed radial derivative
orders. The original derivatives satisfy
\NSi{NS-F-P148-029}{\partial_X^r=\Lambda^r\partial_Y^r}. Real data and uniqueness of the fixed
point give a real solution on the real rectangle. Its analyticity in \NSi{NS-F-P148-030}{X} makes
the scalar profiles smooth at the axis. Together with
\NSi{NS-F-P148-031}{E=\sqrt{2X}\phi/\mathsf C}, this gives a smooth Cartesian field by
Equations~\eqref{eq:4.4} to \eqref{eq:4.5}. The scalar lower bound \eqref{eq:B.11} gives
\NSi{NS-F-P148-032}{\Phi>0} for large \NSi{NS-F-P148-033}{\Lambda}, and hence
\NSi{NS-F-P148-034}{\phi>0}. Division by \NSi{NS-F-P148-035}{\phi} in the first equation of
\eqref{eq:B.15} is therefore valid throughout the real rectangle.
\end{proof}

\subsection{Positive sources and the outer-endpoint alternative.}\label{sec:B.4}
The analytic solution must meet the quantitative conditions needed to begin its continuation toward
the outer profile. We derive a positive angular source \NSi{NS-F-P148-036}{S_q} and a strict shear
inequality \eqref{eq:B.19} at the end of the inner interval from the preceding approximation
\eqref{eq:B.13}. The same estimates will bound the normalized profiles and their parameter
derivatives for use in the joining construction (Proposition~\ref{prop:B.3} and
Lemma~\ref{lem:B.4}).

\begin{proposition}\label{prop:B.3}
Increase \NSi{NS-F-P148-037}{\Lambda} and then \NSi{NS-F-P148-038}{\mathsf C_0(\Lambda)} if necessary.
The normalized azimuthal profile satisfies \NSi{NS-F-P148-039}{\Phi\geq c_0>0} on
\NSi{NS-F-P148-040}{[0,4.1]\times[-1,1]}. On that rectangle its angular source satisfies
\NSFormula{NS-F-P148-041}{display}{B.17}
\begin{equation}
  S_q\geq2.5+.95\Lambda L\chi.
  \tag{B.17}\label{eq:B.17}
\end{equation}
Let \NSi{NS-F-P148-042}{p_1=p_{s,1}=XQ_s/L},
\NSi{NS-F-P148-043}{p_2=p_{s,2}=XN_s/(LE)}, and \NSi{NS-F-P148-044}{n_s=N_s/L}. On
\NSi{NS-F-P148-045}{0<X\leq4.1/\Lambda} one has
\NSFormula{NS-F-P148-046}{display}{B.18}
\begin{equation}
  p_1/X\geq c_1>0,
  \qquad 0<p_1\leq C_1,
  \qquad n_s=-2U_X=Z_*/L+O(\Lambda^{-1}).
  \tag{B.18}\label{eq:B.18}
\end{equation}
\NSPage{149}%
At \NSi{NS-F-P149-001}{X_0=4/\Lambda} there is a fixed \NSi{NS-F-P149-002}{c_{\mathrm{ex}}>0}
such that
\NSFormula{NS-F-P149-003}{display}{B.19}
\begin{equation}
  p_1+\frac{p_2^2}{p_1}>2+c_{\mathrm{ex}}.
  \tag{B.19}\label{eq:B.19}
\end{equation}
The quantities \NSi{NS-F-P149-004}{\Phi,\log\Phi,\mathfrak u,p_1}, and \NSi{NS-F-P149-005}{n_s} have
bounded derivatives in every fixed parameter order, uniformly in subsequent
\NSi{NS-F-P149-006}{\mathsf C\geq\mathsf C_0(\Lambda)}; constants for derivatives in the original radial variable
\NSi{NS-F-P149-007}{X} may depend on \NSi{NS-F-P149-008}{\Lambda}.
\end{proposition}

\begin{proof}
The scalar bound \eqref{eq:B.11} and the approximation \eqref{eq:B.13} give
\NSi{NS-F-P149-009}{\Phi\geq c_0>0} for large \NSi{NS-F-P149-010}{\Lambda}. On a smaller
common parameter neighborhood \NSi{NS-F-P149-011}{\Phi} is also nonzero, by compactness and the
uniform first derivative bound. The parameter derivatives of \NSi{NS-F-P149-012}{\log\Phi} are
consequently uniformly bounded.

The source bound depends on the logarithmic derivatives of \NSi{NS-F-P149-013}{\Phi}. Since
\NSi{NS-F-P149-014}{|H_*\chi'|\leq2\lVert H_*'\rVert_\infty\chi} and
\NSi{NS-F-P149-015}{f_0} is bounded away from zero,
\NSi{NS-F-P149-016}{\mathcal D_X\log f_0(Y\chi)=O(\chi)} and
\NSi{NS-F-P149-017}{H_*\partial_\eta\log f_0(Y\chi)=O(\chi)}. The approximation and positivity
therefore imply for the full solution
\NSFormula{NS-F-P149-018}{display}{B.20}
\begin{equation}
  |\mathcal D_X\log\Phi|+|H_*\partial_\eta\log\Phi|\leq C\chi+C/\Lambda.
  \tag{B.20}\label{eq:B.20}
\end{equation}
Using \eqref{eq:B.14} and \eqref{eq:B.3}, this gives
\NSFormula{NS-F-P149-019}{display}{B.21}
\begin{align}
  -H_c\widetilde\xi_0&=\Lambda L\chi+O_{\sigma_*}(\sqrt\chi),\notag\\
  -H_c(\log\phi)_\eta
    &\geq\Lambda L\chi-C\chi-C_{\sigma_*}\sqrt\chi-C/\Lambda
      \geq.95\Lambda L\chi-C_{\sigma_*}/\Lambda.
  \tag{B.21}\label{eq:B.21}
\end{align}
For example,
\NSi{NS-F-P149-020}{|H_*|/(H_*^2+\sigma_*^2)\leq\sigma_*^{-1}\sqrt\chi} controls the first
error. The second inequality follows by increasing \NSi{NS-F-P149-021}{\Lambda} to absorb
\NSi{NS-F-P149-022}{C\chi} and then applying Young's inequality to the square-root term, allocating
a total of \NSi{NS-F-P149-023}{.05\Lambda L\chi}.

For the source itself use \NSi{NS-F-P149-024}{l=1+\mathcal D_X\log\Phi} and the defining formula
\NSi{NS-F-P149-025}{S_q=-Wl-h(1-2\eta U)-H_c(\log E)_\eta}. By
\NSi{NS-F-P149-026}{j_0\leq.05} and \eqref{eq:B.13}, increasing
\NSi{NS-F-P149-027}{\Lambda} ensures \NSi{NS-F-P149-028}{|U|\leq4.1} and therefore
\NSi{NS-F-P149-029}{|h(1-2\eta U)|\leq10h}. The preceding bounds and
\NSi{NS-F-P149-030}{W-W_*=O(\Lambda^{-1})} imply
\NSFormula{NS-F-P149-031}{display}{none}
\begin{equation*}
  S_q\geq\Lambda L\chi-W_*-10h-C\chi-C_{\sigma_*}\sqrt\chi-C/\Lambda.
\end{equation*}
Since \NSi{NS-F-P149-032}{h\leq10^{-2}} and \NSi{NS-F-P149-033}{-W_*>2.8}, the same absorption
with \NSi{NS-F-P149-034}{\Lambda} large proves \eqref{eq:B.17}. In particular no division by
\NSi{NS-F-P149-035}{\chi} is used near a zero of \NSi{NS-F-P149-036}{H_*}.

The function \NSi{NS-F-P149-037}{Q_s} extends smoothly to \NSi{NS-F-P149-038}{X=0} and has
the positive integral representation
\NSFormula{NS-F-P149-039}{display}{none}
\begin{equation*}
  Q_s(X,\eta)
  =\frac{\displaystyle\int_0^X X'\Phi(\Lambda X',\eta)S_q(X',\eta)\,dX'}
         {X^2\Phi(\Lambda X,\eta)}.
\end{equation*}
Uniform upper and lower bounds for \NSi{NS-F-P149-040}{\Phi} and \eqref{eq:B.17} give
\NSi{NS-F-P149-041}{p_1/X\geq c_1}. Integrating the equations with vanishing leading residual
stress from \NSi{NS-F-P149-042}{X=0} and using regularity at the axis gives the identities
\NSFormula{NS-F-P149-043}{display}{none}
\begin{equation*}
  p_1=a=-2Y\Phi_Y/\Phi,
  \qquad n_s=-2\mathfrak u_Y.
\end{equation*}
The upper bound for \NSi{NS-F-P149-044}{p_1} and the approximation for
\NSi{NS-F-P149-045}{n_s} now follow from \eqref{eq:B.13}. These identities also give the claimed
fixed bounds on \NSi{NS-F-P149-046}{\eta}-derivatives.

It remains to prove the endpoint inequality. The two parameter regions selected by \eqref{eq:B.2}
allow us to use either the azimuthal or the axial contribution. At \NSi{NS-F-P149-047}{Y=4}, first
suppose \NSi{NS-F-P149-048}{\chi>.99}, so \NSi{NS-F-P149-049}{1.98<t=2\chi\leq2}. The
alternating series for \NSi{NS-F-P149-050}{f_0(z)+zf_0'(z)} gives
\NSFormula{NS-F-P149-051}{display}{none}
\begin{equation*}
  f_0(z)+zf_0'(z)\leq1-t+\frac{t^2}{4}-\frac{t^3}{36}+\frac{t^4}{576}<-.18.
\end{equation*}
The polynomial is decreasing on \NSi{NS-F-P149-052}{[1.98,2]}. At
\NSi{NS-F-P149-053}{t=1.98} its value is
\NSi{NS-F-P149-054}{-75535511/400000000<-.188}. The decreasing alternating remainder has the
required negative sign. By \eqref{eq:B.11},
\NSPage{150}%
\NSi{NS-F-P150-001}{-2zf_0'/f_0=2-2(f_0+zf_0')/f_0>2.36}. Equation \eqref{eq:B.13}
therefore gives \NSi{NS-F-P150-002}{p_1>2.3} at the outer endpoint for large
\NSi{NS-F-P150-003}{\Lambda}.

In the complementary region \NSi{NS-F-P150-004}{\chi\leq.99}, \eqref{eq:B.2} gives
\NSi{NS-F-P150-005}{|Z_*|>\delta_*}. Thus \eqref{eq:B.18} implies
\NSi{NS-F-P150-006}{|n_s|\geq\delta_*/2}, after increasing
\NSi{NS-F-P150-007}{\Lambda}. At the outer endpoint,
\NSFormula{NS-F-P150-008}{display}{none}
\begin{equation*}
  |p_2|=\mathsf C\sqrt{2/\Lambda}\frac{|n_s|}{\phi_*\Phi}.
\end{equation*}
For the now fixed \NSi{NS-F-P150-009}{\Lambda}, the denominator has a finite upper bound uniform
in subsequent large \NSi{NS-F-P150-010}{\mathsf C}, whereas \NSi{NS-F-P150-011}{p_1\leq C_1}.
Increasing \NSi{NS-F-P150-012}{\mathsf C_0(\Lambda)} therefore makes
\NSi{NS-F-P150-013}{p_2^2/p_1>2.3} uniformly on this region. Both alternatives imply
\eqref{eq:B.19}, for example with \NSi{NS-F-P150-014}{c_{\mathrm{ex}}=.2}.
\end{proof}

\subsection{A reference continuation satisfying
\NSi{NS-F-P150-015}{p_{s,1}+p_{s,2}^2/p_{s,1}>2}.}\label{sec:B.5}
We next continue the analytic axis profile of Proposition~\ref{prop:B.2} toward the outer profile of
Section~\ref{sec:A}. The continuation must introduce nonzero leading residual stress
\NSi{NS-F-P150-016}{\mathcal T_0} satisfying the cone conditions and preserve the pressure datum
\NSi{NS-F-P150-017}{\Pi_0} of Lemma~\ref{lem:A.5}. In this subsection, we first construct a
reference continuation \NSi{NS-F-P150-018}{\phi_{\mathrm r},U_{\mathrm r}} by gradually setting the logarithmic radial
derivatives of \NSi{NS-F-P150-019}{\phi} and \NSi{NS-F-P150-020}{U} to zero, as prescribed below
in \eqref{eq:B.22}. Bounds on its integrated residual coefficients, proved in Lemma~\ref{lem:B.4},
will control the subsequent change in shear (Proposition~\ref{prop:B.5}).

As in Proposition~\ref{prop:B.3}, \NSi{NS-F-P150-021}{p_1,p_2} denote the components of
\NSi{NS-F-P150-022}{p_s}, and \NSi{NS-F-P150-023}{n_s=N_s/L}. Recall that
\NSi{NS-F-P150-024}{\mathcal D_X=X\partial_X}, equivalently \NSi{NS-F-P150-025}{\partial_y} in a
logarithmic radial coordinate. Bounds in \NSi{NS-F-P150-026}{C_\eta^k} refer only to parameter
derivatives, uniformly for \NSi{NS-F-P150-027}{-1\leq\eta\leq1}.

Let \NSi{NS-F-P150-028}{X_0=4/\Lambda}, \NSi{NS-F-P150-029}{X_{\mathrm b}=100}, and
\NSi{NS-F-P150-030}{X_i=110}. We choose a fixed \NSi{NS-F-P150-031}{\bar t>0} with
\NSi{NS-F-P150-032}{4e^{2\bar t}<4.1}. We use the smooth step \NSi{NS-F-P150-033}{\sigma} of
\eqref{eq:A.5}. For \NSi{NS-F-P150-034}{0<t_1\leq\bar t}, put
\NSi{NS-F-P150-035}{y=\log(X/X_0)}. The reference profile agrees with the analytic axis profile for
\NSi{NS-F-P150-036}{y\leq t_1}; for \NSi{NS-F-P150-037}{t_1<y<2t_1} we prescribe
\NSFormula{NS-F-P150-038}{display}{B.22}
\begin{align}
  \partial_y\log\phi_{\mathrm r}
    &=[1-\sigma((y-t_1)/t_1)]\mathcal D_X\log\phi_{\mathrm{nat}}(X,\eta),\notag\\
  \partial_yU_{\mathrm r}
    &=[1-\sigma((y-t_1)/t_1)]\mathcal D_XU_{\mathrm{nat}}(X,\eta),
  \tag{B.22}\label{eq:B.22}
\end{align}
and then extend both fields constantly in the logarithmic radial coordinate. We compute the
reference pressure and the functions \NSi{NS-F-P150-039}{Q_s,N_s} by radial integration from the
axis, using the datum of Lemma~\ref{lem:A.5}.

\begin{lemma}\label{lem:B.4}
After the parameters of the outer and axis profiles are fixed, one can take
\NSi{NS-F-P150-040}{\Lambda}, then \NSi{NS-F-P150-041}{\mathsf C}, sufficiently large, and then
\NSi{NS-F-P150-042}{t_1} sufficiently small, so that on
\NSi{NS-F-P150-043}{[X_0,X_i]\times[-1,1]}
\NSFormula{NS-F-P150-044}{display}{B.23}
\begin{equation}
  S_{q,\mathrm r}\geq.94L\Lambda\chi+2.4,
  \qquad l_{\mathrm r}\leq1,
  \qquad v_{\mathrm r}:=p_{1,\mathrm r}+\frac{p_{2,\mathrm r}^2}{p_{1,\mathrm r}}>2+c.
  \tag{B.23}\label{eq:B.23}
\end{equation}
Here \NSi{NS-F-P150-045}{p_{1,\mathrm r}>0}, and \NSi{NS-F-P150-046}{p_{1,\mathrm r}>3} for
\NSi{NS-F-P150-047}{X_{\mathrm b}\leq X\leq X_i}. For each fixed \NSi{NS-F-P150-048}{k}, the parameter
norms of \NSi{NS-F-P150-049}{\mathsf C E_{\mathrm r}}, its reciprocal, \NSi{NS-F-P150-050}{p_{1,\mathrm r}}, and
\NSi{NS-F-P150-051}{n_{s,\mathrm r}} on \NSi{NS-F-P150-052}{[X_0,X_i]} have bounds independent of
sufficiently large \NSi{NS-F-P150-053}{\mathsf C} and of \NSi{NS-F-P150-054}{0<t_1\leq\bar t}.
\end{lemma}

\begin{proof}
The analytic axis profile has a positive source and a strict inequality at its outer endpoint by
Propositions~\ref{prop:B.2} and \ref{prop:B.3}. We use its gradient bounds \eqref{eq:B.20} and
\eqref{eq:B.21} and source bound \eqref{eq:B.17} to follow these properties through the cutoff.
Integrating \eqref{eq:B.22} shows that \NSi{NS-F-P150-055}{U_{\mathrm r}} and its radial average stay within
\NSi{NS-F-P150-056}{O(\Lambda^{-1})+o_{t_1}(1)} of \NSi{NS-F-P150-057}{U_*} in each required
parameter norm. Indeed radial averaging is a contraction:
\NSFormula{NS-F-P150-058}{display}{none}
\begin{equation*}
  \sup_{0<X\leq X_i}\lVert\mathcal A_X(U_{\mathrm r})-U_*\rVert_{C_\eta^k}
  \leq\sup_{0\leq X\leq X_i}\lVert U_{\mathrm r}-U_*\rVert_{C_\eta^k}.
\end{equation*}
The estimate is independent of the length of the interval on which the reference profile is
constant. The same slope integration bounds \NSi{NS-F-P150-059}{\log\phi_{\mathrm r}-\log\phi_*} in each
parameter norm,
